\section{Introduction}
\label{sec:intro}

When an agent cannot see everything that matters, it has to remember. The standard test of
whether memory helps is a comparison: train a recurrent agent and a feedforward one on a partially
observable task and see which does better. If the recurrent agent wins, the gap is read as the
value of memory. Small, fast environments such as XLand-MiniGrid \citep{nikulin2024xland} make
such comparisons cheap enough to run on a laptop, and so they are run often.

We started from one such comparison. In a pilot with three seeds on an XLand-MiniGrid pickup task,
an LSTM agent learned the task and a memory-free MLP did not; on held-out maps the LSTM was ahead
by 17 to 42 percentage points on every seed. It looked like a clean memory result. Before building
on it we asked two questions that apply to any result of this kind.

\textbf{How much of a measured gap is noise?} A run's score depends on the training seed, on the
maps it is tested on, and on the randomness of the evaluation episodes themselves. With three
seeds and a single evaluation per checkpoint these sources are mixed, and some popular summaries,
such as a yes/no verdict that a task was ``acquired'', may be dominated by the last one. Prior
work has shown that deep RL comparisons need more seeds than are usually run
\citep{henderson2018matters,colas2018seeds,agarwal2021precipice,patterson2024empirical}; we ask
how many, for which metric, and why, by separating the three sources of variation with the
tools of generalizability theory \citep{cronbach1972dependability,brennan2001generalizability}.

\textbf{Is the gap memory?} The LSTM in our pilot was also four times larger than the MLP, and
it could do something the MLP could not: count the steps of the episode. Episodes are cut at a
fixed time limit and the reward shrinks with time, so knowing the time is useful on its own
\citep{pardo2018timelimits}. A memory advantage is only an advantage of memory once capacity and
time awareness are controlled.

We answer both questions with 188 runs on one laptop CPU, every one listed in a plan before its
results existed, executed as an event-driven chain and re-evaluated from saved weights. Our
contributions are:
\begin{itemize}
  \item \textbf{A variance decomposition of an RL comparison} into training seed, map, seed
  $\times$ map and evaluation noise (60 runs, \cref{sec:results}), with the number of seeds each
  metric needs. About ten paired seeds detect a 10-point gap in final success, but the binary
  acquisition verdict is 76\% evaluation noise and needs about a hundred.
  \item \textbf{An audit of the memory advantage} with parameter-matched controls, a
  recurrence ablation, a short fixed memory and a clock input, tested on a sealed panel of new
  maps opened once (\cref{sec:part2}). A memory-free MLP of the LSTM's size that is told the time
  step comes within 3.3 points of the LSTM; width without the clock adds nothing measurable.
  Most of the apparent memory advantage is an advantage of knowing the time.
  \item \textbf{Evidence for the mechanism}: a linear probe reads the time step from the trained
  LSTM's state, but not from the observation or from an LSTM whose state is erased; and giving
  the time step to that erased LSTM, or to an MLP a quarter of the LSTM's size, brings either to
  within 5 points of the LSTM.
  \item \textbf{Tooling and data}: preregistered plans, an event-driven runner with no silent
  retries, re-evaluation of every checkpoint, sealed panels, and all weights and per-episode
  outcomes. Re-running a seed reproduces its weights bit for bit.
\end{itemize}
